\documentclass[11pt]{article}
\usepackage[utf8]{inputenc}
\usepackage{amsmath,amssymb}
\usepackage[margin=1in]{geometry}
\usepackage{hyperref}
\emergencystretch=3em

\title{The constant-Gaussian fluctuation model: Tonelli identity}
\author{Claude, with Daniel Murfet}
\date{2026-09-07}

\begin{document}
\maketitle
\sloppy

\begin{abstract}
First half of the rem:pop\_vs\_emp dichotomy (Astra~\#12~(b), \#13~(c)). The leading coefficient is
$A_p=\frac{y_{00}}{k_1k_2}J_p(x_{00})$ with the positive Gaussian moment
$J_p(x)=\int_0^\infty s^{p-1}e^{-\beta s^2}e^{\beta sx}ds$. For a Gaussian constant fluctuation
$X\sim N(0,v)$ we prove, in the extended nonnegative reals,
\[
\mathbb E_+[J_p(X)]=\int_0^\infty s^{p-1}e^{-\beta s^2}e^{v(\beta s)^2/2}\,ds
\]
(\texttt{lintegral\_gaussMomentJ\_eq}), by Tonelli and the Gaussian moment generating function.
The right-hand side is finite iff $\beta v<2$ (next unit): the expected leading coefficient can be
infinite although every sample coefficient is finite. Zero \texttt{sorry}/\texttt{axiom}.
\end{abstract}

\section{The things formalised}
{\small
\begin{verbatim}
def gaussMomentJ (β p x : ℝ) : ℝ := logMoment β p 0 (fun s => Real.exp (β * s * x))
theorem gaussMomentJ_pos (hβ) (hp) : 0 < gaussMomentJ β p x
theorem measurable_gaussMoment_integrand (β p) (X : Ω → ℝ) (hX : Measurable X) :
    Measurable fun q : Ω × ℝ => ENNReal.ofReal (q.2^(p-1) * log q.2 ^ 0
      * (exp (-β * q.2^2) * exp (β * q.2 * X q.1)))
theorem lintegral_exp_gaussian (X) (hX : Measurable X) (v : NNReal)
    (hXg : μ.map X = gaussianReal 0 v) (t) :
    ∫⁻ ω, ENNReal.ofReal (exp (t * X ω)) ∂μ = ENNReal.ofReal (exp (v * t ^ 2 / 2))
theorem lintegral_gaussMomentJ_eq (β p) (hβ) (hp) (X) (hX) (v) (hXg) :
    ∫⁻ ω, ENNReal.ofReal (gaussMomentJ β p (X ω)) ∂μ
      = ∫⁻ s in Ioi 0, ENNReal.ofReal (s^(p-1) * exp (-β * s^2) * exp (v * (β * s)^2 / 2))
\end{verbatim}
}

\section{Mathematics}
Each $J_p(X(\omega))$ is a Bochner integral of a positive integrable function, hence equal (after
\texttt{ofReal}) to the corresponding Lebesgue integral. The joint integrand is measurable, so
Tonelli swaps the order; the inner integral over $\omega$ is the exponential moment
$\mathbb E[e^{\beta sX}]=e^{v(\beta s)^2/2}$ (Mathlib's \texttt{mgf\_gaussianReal}, transported
from the pushforward measure by \texttt{integrable\_map\_measure}). Isolating the positive moment
$J_p$ from the prefactor $y_{00}/(k_1k_2)$, as Astra~\#12 advised, keeps the sign cases out of
the core identity.

\section{Paper-facing meaning}
rem:pop\_vs\_emp states that the expected empirical coefficient differs from the population one
because the coefficients depend nonlinearly on the fluctuation; the present identity quantifies this
for the simplest random fluctuation (a Gaussian constant): the expectation replaces $e^{-\beta s^2}$ by
$e^{-(\beta-\beta^2v/2)s^2}$, so $\mathbb E[J_p(X)]=J_p(0)(1-\beta v/2)^{-p/2}>J_p(0)$ when
$\beta v<2$ and $=+\infty$ when $\beta v\ge2$ (next unit). As Astra~\#12 stressed, this concerns
the expected asymptotic coefficient, not the fixed-$n$ chart integral on a compact chart, whose
expectation stays finite.

\section{Lean notes}
\begin{itemize}
\item \texttt{integrable\_map\_measure hg hX.aemeasurable} transports integrability along
\texttt{μ.map X = gaussianReal 0 v}; the a.e.-strong-measurability argument must be given as a
separate \texttt{have} (an inline dot-chain mis-associates).
\item \texttt{lintegral\_lintegral\_swap} needs \texttt{AEMeasurable (uncurry f)} on the product;
\texttt{Measurable.pow\_const} covers real powers of the coordinate.
\item \texttt{ENNReal.ofReal\_mul} (with the nonnegativity of the left factor) and
\texttt{lintegral\_const\_mul} pull the $s$-dependent factor out of the $\omega$-integral.
\end{itemize}

\end{document}
